\documentclass[11pt]{amsart}

\usepackage{amsmath,amssymb,amsthm}
\usepackage[margin=2.6cm]{geometry}
\usepackage{booktabs}
\usepackage{longtable}
\usepackage{seqsplit}
\usepackage[hidelinks]{hyperref}

\newtheorem{theorem}{Theorem}[section]
\newtheorem{lemma}[theorem]{Lemma}
\newtheorem{proposition}[theorem]{Proposition}
\newtheorem{corollary}[theorem]{Corollary}
\theoremstyle{definition}
\newtheorem{definition}[theorem]{Definition}
\theoremstyle{remark}
\newtheorem{remark}[theorem]{Remark}

\newcommand{\lean}[1]{\texttt{\seqsplit{#1}}}
\newcommand{\eqb}{\equiv}
\newcommand{\conj}{\sim}
\newcommand{\rot}{\approx}
\DeclareMathOperator{\len}{len}

\title[J08: Coxeter conjugacy of positive words]{J08: centrality of the full twist and conjugacy
of Coxeter words\\ in the positive braid monoid}
\date{2 November 2026}

\begin{document}

\begin{abstract}
We give the rewriting sequences behind job J08 (record M408) of the dependency audit. On
$m$ strands, with $c=\sigma_0\sigma_1\cdots\sigma_{m-2}$ and $\Delta^2=c^m$, we show that
$\Delta^2\sigma_i$ and $\sigma_i\Delta^2$ are related by far commutations and positive braid
relations applied inside words, for every generator $\sigma_i$. We then show that for $m\ge2$ every
word $\gamma$ containing each generator exactly once satisfies $\Delta^2\gamma\sim\Delta^2c$, where
$\sim$ is the relation \lean{PositiveWordConj}. For $m\le1$ the second statement is false, because
the only word is the empty word and \lean{PositiveWordConj} relates no word to the empty word. All
proofs are elementary and self-contained.
\end{abstract}

\maketitle

\section{Words and relations}\label{sec:setting}

Fix an integer $m\ge0$ and put $n:=\max(m-1,0)$, the value of the truncated difference $m-1$ in
the Lean code. The \emph{alphabet} is $\{0,1,\dots,n-1\}$, and the letter $i$ stands for the
positive generator $\sigma_i$ of the braid group on $m$ strands. A \emph{word} is a finite sequence of
letters (the Lean type \lean{PositiveAdjacentWord m}, a list of elements of \lean{Fin (m - 1)}).
Juxtaposition $uv$ denotes concatenation, $\varnothing$ is the empty word and $\len(u)$ is the length
of $u$. Two letters $i,j$ are \emph{far} if $|i-j|\ge2$. For a letter $x$ and a word $w$ we write
$x\perp w$ if $x$ is far from every letter of $w$.

\begin{definition}\label{def:moves}
An \emph{elementary braid move} replaces a word $u\,i\,j\,v$ with $i,j$ far by $u\,j\,i\,v$, or a
word $u\,i\,(i+1)\,i\,v$ by $u\,(i+1)\,i\,(i+1)\,v$, where $u,v$ are arbitrary words and
$i+1\le n-1$ in the second case. A \emph{rotation} replaces a word $u\,i$ by $i\,u$, where $i$ is a
letter.
\end{definition}

We use three relations on words.
\begin{enumerate}
\item \emph{Positive braid equivalence} $u\eqb v$ is the equivalence relation generated by elementary
braid moves (Lean: \lean{PositiveBraidEquiv}, the closure \lean{Relation.EqvGen} of
\lean{PositiveBraidStep}).
\item \emph{Positive word conjugacy} $u\conj v$ is the relation \lean{PositiveWordConj} of
\lean{PositiveWordNormalForms.lean}: the smallest relation containing all pairs $(u\,i,\ i\,u)$ given
by rotations and all pairs $(u',u'')$ given by elementary braid moves, and closed under symmetry and
transitivity. Reflexivity is not part of the definition.
\item \emph{Far rotation equivalence} $u\rot v$ is the equivalence relation generated by elementary
braid moves of the first kind (far commutations) and by rotations (Lean:
\lean{FarRotationEquiv}).
\end{enumerate}

\begin{lemma}\label{lem:basic}
Let $u,v$ be words.
\begin{enumerate}
\item If $u\eqb v$, then $a\,u\,b\eqb a\,v\,b$ for all words $a,b$.
\item If $u\conj v$, then $\len(u)=\len(v)\ge1$. In particular $\varnothing\conj v$ holds for no word $v$.
\item If $u\ne\varnothing$, then $u\conj u$.
\item If $u\eqb v$ and $u\ne\varnothing$, then $u\conj v$.
\end{enumerate}
\end{lemma}

\begin{proof}
(1) An elementary braid move from $u$ to $v$ remains an elementary braid move from $aub$ to $avb$,
after the outer words of Definition~\ref{def:moves} are enlarged by $a$ and $b$. The set of pairs
$(u,v)$ with $aub\eqb avb$ for all $a,b$ is therefore an equivalence relation containing the
elementary braid moves, and it contains $\eqb$.

(2) Rotations and elementary braid moves preserve the length, and each of them relates two nonempty
words. Symmetry and transitivity preserve both properties.

(3) Write $u=u'x$ with a letter $x$. The rotation gives $u\conj xu'$, symmetry gives $xu'\conj u$,
and transitivity gives $u\conj u$.

(4) The set of pairs $(u,v)$ with $\len(u)=\len(v)$ and with $u\conj v$ whenever $u\ne\varnothing$ is
an equivalence relation by (3) and contains the elementary braid moves. Hence it contains $\eqb$.
\end{proof}

\begin{remark}\label{rem:m2}
Rotations enter the proof of Corollary~\ref{cor:central-conj} below only through
Lemma~\ref{lem:basic}(3), which provides $u\conj u$ for the reflexivity steps of $\eqb$. Every other
step is an elementary braid move inside a word. If $m=2$, the alphabet is $\{0\}$, no elementary
braid move exists, $\eqb$ is equality, and the relation $u\conj u$ cannot be obtained without a
rotation.
\end{remark}

\begin{lemma}\label{lem:pass}
Let $x$ be a letter and $w$ a word with $x\perp w$. Then $w\,x\eqb x\,w$ and $w\,x\rot x\,w$, and the
same equivalences hold after adding arbitrary words on the left and on the right.
\end{lemma}

\begin{proof}
Induction on $\len(w)$. For $w=y\,w'$ the induction hypothesis gives $y\,w'\,x\eqb y\,x\,w'$, and
the far commutation of $y$ and $x$ gives $y\,x\,w'\eqb x\,y\,w'$. The resulting chain from $w\,x$ to
$x\,w$ consists of far commutations. A far commutation remains a far commutation after arbitrary words
are added on the left and on the right, so the chain also proves $a\,w\,x\,b\eqb a\,x\,w\,b$ and
$a\,w\,x\,b\rot a\,x\,w\,b$ for all words $a,b$.
\end{proof}

\section{Centrality of the full twist}\label{sec:central}

Assume $n\ge1$, that is, $m\ge2$. Let $c:=0\,1\,\cdots\,(n-1)$ be the Coxeter word
(\lean{coxeterWord m}), let $c^k$ be its $k$-th power (\lean{torusWord m k}), and let
$\Delta^2:=c^m=c^{n+1}$ (\lean{fullTwistWord m}). For a word $w$ whose letters are at most $n-2$,
let $w^+$ be the word obtained by adding $1$ to every letter. Put
\[
\hat c:=0\,1\cdots(n-2),\qquad \tilde c:=1\,2\cdots(n-1),
\]
so that $c=\hat c\,(n-1)=0\,\tilde c$ and $\hat c^{\,+}=\tilde c$. For $n=1$ both words are empty.

\begin{lemma}\label{lem:shift}
\begin{enumerate}
\item For every letter $i\le n-2$, $c\,i\eqb(i+1)\,c$.
\item For every word $w$ whose letters are at most $n-2$, $c\,w\eqb w^+\,c$.
\item For all $k\ge0$ and every letter $i$ with $i+k\le n-1$, $c^k\,i\eqb(i+k)\,c^k$.
\end{enumerate}
\end{lemma}

\begin{proof}
(1) Write $c=A\,i\,(i+1)\,B$ with $A=0\cdots(i-1)$ and $B=(i+2)\cdots(n-1)$. Every letter of $B$ is
at least $i+2$, so $i\perp B$, and every letter of $A$ is at most $i-1$, so $(i+1)\perp A$. By
Lemma~\ref{lem:pass}, the braid move, and Lemma~\ref{lem:pass} again,
\[
c\,i=A\,i\,(i+1)\,B\,i\ \eqb\ A\,i\,(i+1)\,i\,B\ \eqb\ A\,(i+1)\,i\,(i+1)\,B\ \eqb\
(i+1)\,A\,i\,(i+1)\,B=(i+1)\,c .
\]

(2) Induction on $\len(w)$. For $w=x\,w'$, part (1) and the induction hypothesis give
$c\,x\,w'\eqb(x+1)\,c\,w'\eqb(x+1)\,w'^{+}\,c=w^+c$, by Lemma~\ref{lem:basic}(1).

(3) Induction on $k$, for all $i$ at once. The case $k=0$ is trivial. If $i+k+1\le n-1$, then
$i\le n-2$, and part (1) and the induction hypothesis applied to the letter $i+1$ give
\[
c^{k+1}\,i=c^k\,c\,i\ \eqb\ c^k\,(i+1)\,c\ \eqb\ (i+1+k)\,c^k\,c=(i+k+1)\,c^{k+1}. \qedhere
\]
\end{proof}

\begin{lemma}\label{lem:wrap}
$c^2\,(n-1)\eqb 0\,c^2$.
\end{lemma}

\begin{proof}
The letters of $\hat c$ are at most $n-2$, so Lemma~\ref{lem:shift}(2) gives $c\,\hat c\eqb\tilde c\,c$.
Hence, by Lemma~\ref{lem:basic}(1),
\[
0\,c^2=0\,c\,\hat c\,(n-1)\ \eqb\ 0\,\tilde c\,c\,(n-1)=c\,c\,(n-1)=c^2\,(n-1). \qedhere
\]
\end{proof}

\begin{theorem}\label{thm:central}
For every letter $i$, $\Delta^2\,i\eqb i\,\Delta^2$. Consequently
$u\,x\,\Delta^2\,v\eqb u\,\Delta^2\,x\,v$ for all words $u,x,v$.
\end{theorem}

\begin{proof}
Let $k:=n-1-i$, so that $\Delta^2=c^{n+1}=c^i\,c^2\,c^k$. Lemma~\ref{lem:shift}(3) with the letter
$i$ and exponent $k$, Lemma~\ref{lem:wrap}, and Lemma~\ref{lem:shift}(3) with the letter $0$ and
exponent $i$ give
\[
\Delta^2\,i=c^i\,c^2\,c^k\,i\ \eqb\ c^i\,c^2\,(n-1)\,c^k\ \eqb\ c^i\,0\,c^2\,c^k\ \eqb\
i\,c^i\,c^2\,c^k=i\,\Delta^2 ,
\]
where each step uses Lemma~\ref{lem:basic}(1). The second statement follows by induction on
$\len(x)$: for $x=y\,x'$,
$u\,y\,x'\,\Delta^2\,v\eqb u\,y\,\Delta^2\,x'\,v\eqb u\,\Delta^2\,y\,x'\,v$.
\end{proof}

\begin{corollary}\label{cor:central-conj}
Let $m\ge0$. For all words $u,v$ and every letter $i$,
$u\,i\,\Delta^2\,v\conj u\,\Delta^2\,i\,v$. If $m\ge2$, then $u\,x\,\Delta^2\,v\conj u\,\Delta^2\,x\,v$
for all words $u,x,v$.
\end{corollary}

\begin{proof}
If a letter exists, then $m\ge2$, and the first word is nonempty. Theorem~\ref{thm:central} and
Lemma~\ref{lem:basic}(4) give the first statement. If $m\ge2$, then $\len(\Delta^2)=m(m-1)\ge2$, and
the second statement follows in the same way. If $m\le1$, no letter exists and the first statement
is vacuous.
\end{proof}

\section{Conjugacy of Coxeter words}\label{sec:coxeter}

A \emph{Coxeter word} is a word containing every letter $0,\dots,n-1$ exactly once, that is, a
permutation of the list $c$. For $0\le k\le n$ let $s_k:=0\,1\cdots(k-1)$, the first $k$ letters of
$c$ (Lean: \lean{(List.finRange (m - 1)).take k}), so that $s_n=c$ and $s_{k+1}=s_k\,k$.

\begin{remark}\label{rem:humphreys}
The audit cites Humphreys, \emph{Reflection Groups and Coxeter Groups}, \S3.16, for the conjugacy of
Coxeter elements in a Coxeter group whose Coxeter graph is a tree. That reference was not consulted
for this note. Proposition~\ref{prop:coxeter} below is proved directly for the path
$0-1-\cdots-(n-1)$, by induction on the number of letters instead of by the source and sink argument.
\end{remark}

\begin{lemma}\label{lem:rot}
\begin{enumerate}
\item If $u\rot v$, then $v$ is a permutation of $u$.
\item For all words $\alpha,\beta$, $\alpha\,\beta\rot\beta\,\alpha$.
\end{enumerate}
\end{lemma}

\begin{proof}
(1) Far commutations and rotations permute letters. (2) Induction on $\len(\beta)$, for all
$\alpha$ at once. For $\beta=b\,\beta'$, the induction hypothesis applied to $\alpha b$ and the
inverse of a rotation give
$\alpha\,b\,\beta'\rot\beta'\,\alpha\,b\rot b\,\beta'\,\alpha$.
\end{proof}

\begin{lemma}\label{lem:lift}
Let $t\le n-1$ be a letter and let $\delta,\delta'$ be words with $\delta\rot\delta'$. Assume that
the letters of $\delta$ are pairwise distinct and smaller than $t$. Then $\delta\,t\rot\delta'\,t$.
\end{lemma}

\begin{proof}
By Lemma~\ref{lem:rot}(1), every word in a chain of moves from $\delta$ to $\delta'$ has the two
properties assumed for $\delta$. It is therefore enough to treat the case in which $\delta'$ is
obtained from $\delta$ by one far commutation or one rotation. A step of the chain in the opposite
direction is handled by this case applied to the reversed step, whose initial word has the same two
properties, followed by symmetry of $\rot$. A far commutation $\delta=u\,i\,j\,v\to u\,j\,i\,v=\delta'$ gives the far
commutation $u\,i\,j\,v\,t\to u\,j\,i\,v\,t$. Let $\delta=u\,i$ and $\delta'=i\,u$. If $i\le t-2$, then
$i$ and $t$ are far and
\[
u\,i\,t\ \rot\ u\,t\,i\ \rot\ i\,u\,t .
\]
Otherwise $i=t-1$. The letters of $u$ are smaller than $t$ and different from $i$, so they are at most
$t-2$ and $t\perp u$. Two rotations and Lemma~\ref{lem:pass} give
\[
u\,i\,t\ \rot\ t\,u\,i\ \rot\ i\,t\,u\ \rot\ i\,u\,t . \qedhere
\]
\end{proof}

\begin{proposition}\label{prop:coxeter}
For $0\le k\le n$, every permutation $\gamma$ of $s_k$ satisfies $\gamma\rot s_k$. In particular
every Coxeter word $\gamma$ satisfies $\gamma\rot c$.
\end{proposition}

\begin{proof}
Induction on $k$. For $k=0$, $\gamma=\varnothing=s_0$. Let $k+1\le n$ and let $\gamma$ be a
permutation of $s_{k+1}=s_k\,k$. Write $\gamma=\alpha\,k\,\beta$. Lemma~\ref{lem:rot}(2) applied to
$\alpha k$ and $\beta$ gives $\gamma\rot\beta\,\alpha\,k$. By Lemma~\ref{lem:rot}(1), the word
$\delta:=\beta\,\alpha$ is a permutation of $s_k$, and the induction hypothesis gives $\delta\rot s_k$.
The letters of $\delta$ are pairwise distinct and smaller than $k$, so Lemma~\ref{lem:lift} with
$t=k$ gives $\delta\,k\rot s_k\,k=s_{k+1}$.
\end{proof}

\begin{theorem}\label{thm:coxeter-conj}
Let $m\ge2$ and let $\gamma$ be a Coxeter word. Then $\Delta^2\gamma\conj\Delta^2c$. If $m\le1$,
then $\Delta^2\gamma\conj\Delta^2c$ is false.
\end{theorem}

\begin{proof}
Let $m\ge2$. We show that $\gamma\rot\gamma'$ implies $\Delta^2\gamma\conj\Delta^2\gamma'$. The set of
pairs $(\gamma,\gamma')$ with $\Delta^2\gamma\conj\Delta^2\gamma'$ is symmetric and transitive, and it
is reflexive by Lemma~\ref{lem:basic}(3), because $\Delta^2\ne\varnothing$. A far commutation
$\gamma=u\,i\,j\,v\to u\,j\,i\,v$ gives the far commutation
$\Delta^2u\,i\,j\,v\to\Delta^2u\,j\,i\,v$. For a rotation $u\,i\to i\,u$, a rotation and
Corollary~\ref{cor:central-conj} give
\[
\Delta^2\,u\,i\ \conj\ i\,\Delta^2\,u\ \conj\ \Delta^2\,i\,u .
\]
Proposition~\ref{prop:coxeter} gives $\gamma\rot c$, hence $\Delta^2\gamma\conj\Delta^2c$.

If $m\le1$, the alphabet is empty, so $\gamma=c=\Delta^2=\varnothing$, and
Lemma~\ref{lem:basic}(2) shows that $\varnothing\conj\varnothing$ is false.
\end{proof}

\section{Invariants of positive word conjugacy}\label{sec:invariants}

For a word $u=x_1\cdots x_r$ let $\pi(u):=\tau_{x_1}\circ\cdots\circ\tau_{x_r}$, where $\tau_i$ is the
transposition of the strands $i$ and $i+1$ (Lean: \lean{PositiveAdjacentWord.permutation}). Let
$L(u)$ be the set of letters of $u$ (Lean: \lean{u.toSigned.generatorIndices}), and let $\kappa(u)$ be
the number of orbits of the cyclic group generated by $\pi(u)$ on $\{0,\dots,m-1\}$, fixed points
included (Lean: \lean{u.toSigned.cycleCount}).

\begin{proposition}\label{prop:invariants}
If $u\conj v$, then $\len(u)=\len(v)$, $L(u)=L(v)$, the permutations $\pi(u)$ and $\pi(v)$ are
conjugate, and $\kappa(u)=\kappa(v)$. If $u\eqb v$, then $\pi(u)=\pi(v)$.
\end{proposition}

\begin{proof}
It suffices to check one rotation and one elementary braid move, since all four properties define
equivalence relations. A rotation permutes letters, and
$\pi(i\,u)=\tau_i\,\pi(u)=\tau_i\,(\pi(u)\,\tau_i)\,\tau_i^{-1}=\tau_i\,\pi(u\,i)\,\tau_i^{-1}$. A far
commutation permutes letters, and $\tau_i\tau_j=\tau_j\tau_i$ for far $i,j$ because the two
transpositions have disjoint supports. A braid move replaces $i\,(i+1)\,i$ by $(i+1)\,i\,(i+1)$,
which has the same set of letters, and both products $\tau_i\tau_{i+1}\tau_i$ and
$\tau_{i+1}\tau_i\tau_{i+1}$ equal the transposition of $i$ and $i+2$. Finally, if $\pi'=g\pi g^{-1}$,
then $g$ maps each orbit of $\pi$ onto an orbit of $\pi'$, so $\kappa$ is a conjugacy invariant.
\end{proof}

\section{Correspondence with the Lean declarations}\label{sec:lean}

The modules \lean{OnLeConjecture.KnotTheory.PositiveWordBraidEquiv} and
\lean{OnLeConjecture.KnotTheory.PositiveWordCoxeterConjugacy} formalize the statements above.

\begin{center}
\begin{longtable}{p{0.33\textwidth}p{0.57\textwidth}}
\toprule
Statement & Lean declaration\\
\midrule
Lemma~\ref{lem:basic}(1) & \lean{PositiveBraidEquiv.append\_left}, \lean{PositiveBraidEquiv.append\_right}, \lean{PositiveBraidEquiv.append}, \lean{PositiveBraidEquiv.append\_middle}\\
Lemma~\ref{lem:basic}(2) & \lean{PositiveWordConj.length\_eq}, \lean{PositiveWordConj.ne\_nil\_left}, \lean{not\_positiveWordConj\_nil\_left}\\
Lemma~\ref{lem:basic}(3) & \lean{PositiveWordConj.refl\_of\_ne\_nil}\\
Lemma~\ref{lem:basic}(4) & \lean{PositiveBraidEquiv.positiveWordConj}, \lean{PositiveBraidEquiv.positiveWordConj\_append\_middle}\\
Lemma~\ref{lem:pass} & \lean{PositiveBraidEquiv.append\_singleton\_of\_far}, \lean{FarRotationEquiv.pass\_far}\\
Lemma~\ref{lem:shift} & \lean{coxeterWord\_append\_singleton}, \lean{coxeterWord\_append\_shift}, \lean{torusWord\_append\_singleton}\\
Lemma~\ref{lem:wrap} & \lean{coxeterWord\_wrap}\\
Theorem~\ref{thm:central} & \lean{torusWord\_strandCount\_append\_singleton}, \lean{fullTwistWord\_append\_singleton}, \lean{fullTwist\_central\_braidEquiv}, \lean{fullTwist\_commute\_braidEquiv}\\
Corollary~\ref{cor:central-conj} & \lean{fullTwist\_central}, \lean{fullTwist\_commute}\\
Lemma~\ref{lem:rot} & \lean{FarRotationEquiv.perm}, \lean{FarRotationEquiv.append\_comm}\\
Lemma~\ref{lem:lift} & \lean{FarRotationStep.append\_singleton\_lift}, \lean{FarRotationEquiv.append\_singleton\_lift}\\
Proposition~\ref{prop:coxeter} & \lean{farRotationEquiv\_take\_finRange}, \lean{farRotationEquiv\_coxeterWord}\\
Theorem~\ref{thm:coxeter-conj} & \lean{FarRotationEquiv.fullTwist\_positiveWordConj}, \lean{fullTwist\_coxeter\_conj}, \lean{fullTwist\_coxeter\_conj\_of\_perm}, \lean{not\_positiveWordConj\_of\_le\_one}, \lean{fullTwist\_coxeter\_conj\_iff}\\
Proposition~\ref{prop:invariants} & \lean{PositiveWordConj.length\_eq}, \lean{PositiveWordConj.generatorIndices\_eq}, \lean{PositiveWordConj.isConj\_permutation}, \lean{PositiveWordConj.cycleCount\_eq}, \lean{PositiveBraidEquiv.permutation\_eq}, \lean{SignedAdjacentWord.cycleCount\_eq\_of\_isConj}\\
\bottomrule
\end{longtable}
\end{center}

\begin{remark}\label{rem:deviation}
The audit proposes \lean{fullTwist\_coxeter\_conj} without a hypothesis on $m$. By
Theorem~\ref{thm:coxeter-conj} that statement is false for $m\in\{0,1\}$. The Lean theorem
\lean{fullTwist\_coxeter\_conj} therefore carries the hypothesis $2\le m$ as its last argument, and
\lean{fullTwist\_coxeter\_conj\_iff} records that the conclusion holds exactly when $2\le m$. The
proposed \lean{fullTwist\_central} is proved without change.
\end{remark}

\end{document}
